\chapter{二重级数与无穷乘积}
\label{chap:double-series-infinite-products}

一列部分和只有一个趋向无穷的指标。二重级数有矩形扩大、先按行求和、
先按列求和以及把所有项排成一列等多种极限过程；缺少非负性或绝对收敛时，
它们可能给出不同结果。无穷乘积也必须通过有限部分乘积定义，并且零极限
与非零极限在理论中承担不同角色。

\section{二重数列、矩形部分和与 Pringsheim 收敛}

设 \((a_{mn})_{m,n\ge1}\) 为二重数列，定义矩形部分和
\[
 S_{MN}=\sum_{m=1}^M\sum_{n=1}^N a_{mn}.
\]

\begin{definition}
若存在 \(s\in\R\)，使对任意 \(\varepsilon>0\)，存在
\(N_0\)，只要 \(M,N\ge N_0\)，就有
\[
 \abs{S_{MN}-s}<\varepsilon,
\]
则称二重级数 \(\sum_{m,n\ge1}a_{mn}\) 按
\term{Pringsheim 意义}收敛于 \(s\)。
\end{definition}

两个边长必须同时趋于无穷，但速度可以不同。只检查方阵
\(S_{NN}\) 不够，因为长而窄的矩形可能保留不同的尾部贡献。

\begin{theorem}[矩形部分和的 Cauchy 准则]
\label{thm:double-series-cauchy}
Pringsheim 收敛等价于：对任意 \(\varepsilon>0\)，存在 \(N_0\)，使
\(M,N,P,Q\ge N_0\) 时
\[
 \abs{S_{MN}-S_{PQ}}<\varepsilon.
\]
\end{theorem}

\begin{proof}
必要性由
\[
 \abs{S_{MN}-S_{PQ}}
 \le\abs{S_{MN}-s}+\abs{S_{PQ}-s}
\]
得到。反之，对每个 \(k\) 取 \(N_k\)，使四个指标都不小于 \(N_k\) 时
两矩形和之差小于 \(2^{-k}\)，并令 \(n_k\ge\max\{N_1,\ldots,N_k,k\}\)
严格递增。对充分大的 \(j,k\)，
\[
 \abs{S_{n_jn_j}-S_{n_kn_k}}
\]
任意小，所以对角序列 \((S_{n_kn_k})\) 是实 Cauchy 列，设其极限为
\(s\)。若 \(M,N\ge N_k\)，先把 \(S_{MN}\) 与一个充分靠后的
\(S_{n_jn_j}\) 比较，再令 \(j\to\infty\)，便得
\(\abs{S_{MN}-s}\le2^{-k}\)。因此矩形部分和按 Pringsheim 意义收敛。
\end{proof}

Pringsheim 收敛只控制以 \((1,1)\) 为左下角的矩形部分和。若还要任意远处
小矩形的和一致变小，需要更强的正则收敛条件；绝对收敛会自动提供这种控制。

\begin{definition}[正则收敛]
\label{def:double-series-regular}
若对任意 \(\varepsilon>0\)，存在 \(N\)，使任意有限矩形
\[
 R=\{m_0,\ldots,m_1\}\times\{n_0,\ldots,n_1\}
\]
只要 \(m_0\ge N\) 或 \(n_0\ge N\)，就有
\[
 \abs{\sum_{(m,n)\in R}a_{mn}}<\varepsilon,
\]
则称二重级数正则收敛。
\end{definition}

正则收敛要求矩形只要沿一个方向远离原点，其和就很小；这正是逐行、逐列
极限与矩形极限能够协调所需的尾部控制。这里只把它作为换序问题的参照，
后续结论主要在更容易核验的绝对收敛条件下陈述。

\section{逐行和、逐列和及换序问题}

若每行级数收敛，可先形成
\[
 r_m=\sum_{n=1}^\infty a_{mn},
\]
再考察 \(\sum_m r_m\)。逐列迭代和类似。这两种极限不必相等，甚至不必
与矩形极限同时存在。

\begin{counterexample}[两种迭代和不同]
\label{ex:double-series-order-counterexample}
令
\[
 a_{m,n}=
 \begin{cases}
 1,&n=m,\\
 -1,&n=m+1,\\
 0,&\text{其他}.
 \end{cases}
\]
每行和为 \(0\)，所以先行后列的和为 \(0\)。第一列和为 \(1\)，其余每列
和为 \(0\)，所以先列后行的和为 \(1\)。此外
\[
 S_{M,M}=1,\qquad S_{M,M+1}=0,
\]
故矩形极限不存在。
\end{counterexample}

这个例子中的每一行、每一列只有有限个非零项，问题并非单行求和本身，
而是无限多行之间缺少统一尾部控制。

\section{非负项与绝对收敛情形}

\begin{theorem}[非负二重级数]
\label{thm:double-series-nonnegative}
若 \(a_{mn}\ge0\)，则扩充实数意义下
\[
 \sup_{M,N}S_{MN}
 =\sum_{m=1}^\infty\left(\sum_{n=1}^\infty a_{mn}\right)
 =\sum_{n=1}^\infty\left(\sum_{m=1}^\infty a_{mn}\right).
\]
因此任一边有限时，三者都有限且相等。
\end{theorem}

\begin{proof}
固定 \(M\)。随 \(N\) 增大，
\[
 S_{MN}\uparrow\sum_{m=1}^M r_m,
\qquad
 r_m=\sum_{n=1}^\infty a_{mn}.
\]
再令 \(M\to\infty\)，得到
\[
 \sum_m r_m=\sup_M\sup_N S_{MN}=\sup_{M,N}S_{MN}.
\]
交换 \(m,n\) 得到列版本。
\end{proof}

\begin{definition}
若
\[
 \sup_{M,N}\sum_{m=1}^M\sum_{n=1}^N\abs{a_{mn}}<\infty,
\]
则称二重级数\term{绝对收敛}。
\end{definition}

由非负二重级数定理，绝对收敛还等价于
\[
 \sum_{m=1}^{\infty}\sum_{n=1}^{\infty}\abs{a_{mn}}<\infty,
\]
也等价于所有有限指标集上的绝对值和具有统一上界。这三个表述分别适合
矩形极限、逐行求和和任意枚举。

\begin{theorem}[绝对收敛下的换序]
\label{thm:double-series-absolute-fubini}
若二重级数绝对收敛，则每行、每列都绝对收敛，两个迭代和与矩形和都存在
且相等。把 \(\Npos^2\) 按任意双射排列为一列后，所得级数也绝对收敛到
同一总和。
\end{theorem}

\begin{proof}
对非负二重数列 \(\abs{a_{mn}}\) 使用上一条定理，得到所有行、列绝对值
之和的总和有限。写
\[
 a_{mn}=a_{mn}^+-a_{mn}^-,
\]
分别对两个非负二重级数换序，即得各求和方式相同。

任意枚举后的绝对值部分和不超过全部绝对值总和，故绝对收敛。给定
\(\varepsilon>0\)，选有限矩形使其外绝对值总和小于 \(\varepsilon\)；
枚举最终包含该矩形全部项，此后与矩形总和的差只来自矩形外，故极限相同。
\end{proof}

\section{Cauchy 乘积、Mertens 定理与条件收敛反例}

从 \(n=0\) 开始写两个级数
\(\sum a_n,\sum b_n\)，定义其 \term{Cauchy 乘积}为
\[
 \sum_{n=0}^\infty c_n,\qquad
 c_n=\sum_{k=0}^n a_kb_{n-k}.
\]
每个 \(c_n\) 沿二重数组的一条反对角线求和。

\begin{theorem}[Mertens 定理]
\label{thm:mertens-cauchy-product}
若 \(\sum a_n\) 绝对收敛、\(\sum b_n\) 收敛，则 Cauchy 乘积收敛，且
\[
 \sum_{n=0}^\infty c_n
 =\left(\sum_{n=0}^\infty a_n\right)
  \left(\sum_{n=0}^\infty b_n\right).
\]
两级数的角色可以交换。
\end{theorem}

\begin{proof}
记 \(A=\sum a_n\)，\(B_N=\sum_{j=0}^N b_j\)，则
\[
 \sum_{n=0}^N c_n
 =\sum_{k=0}^N a_kB_{N-k}.
\]
序列 \((B_N)\) 有界，设 \(\abs{B_N}\le M\)。给定
\(\varepsilon>0\)，先取 \(K\) 使
\[
 \sum_{k>K}\abs{a_k}<\frac{\varepsilon}{3(M+\abs B+1)}.
\]
再取 \(N\) 足够大，使
\[
 \abs{\sum_{k=0}^K a_k(B_{N-k}-B)}<\frac{\varepsilon}{3}.
\]
于是
\[
\begin{aligned}
 \abs{\sum_{n=0}^N c_n-AB}
 &\le
 \abs{\sum_{k=0}^K a_k(B_{N-k}-B)}
 +\abs{\sum_{k=K+1}^N a_kB_{N-k}}\\
 &\qquad
 +\abs{B}\abs{\sum_{k=K+1}^{\infty}a_k}\\
 &<\frac{\varepsilon}{3}
 +(M+\abs B)\sum_{k>K}\abs{a_k}
 <\varepsilon.
\end{aligned}
\]
因此 Cauchy 乘积的部分和趋于 \(AB\)。
\end{proof}

\begin{corollary}
若两个级数都绝对收敛，则其 Cauchy 乘积绝对收敛，且和为两和之积。
\end{corollary}

\begin{proof}
反对角线上的项只有有限个，故三角不等式与非负二重级数的换序给出
\[
 \sum_{n=0}^\infty\abs{c_n}
 \le \sum_{n=0}^\infty\sum_{k=0}^n\abs{a_k}\abs{b_{n-k}}
 =\left(\sum_{k=0}^\infty\abs{a_k}\right)
  \left(\sum_{j=0}^\infty\abs{b_j}\right)<\infty.
\]
和的公式再由 Mertens 定理得到。
\end{proof}

\begin{counterexample}
令
\[
 a_n=b_n=\frac{(-1)^n}{\sqrt{n+1}}\qquad(n\ge0).
\]
两个级数都由 Leibniz 判别收敛。其 Cauchy 乘积通项为
\[
 c_n=(-1)^n\sum_{k=0}^n
 \frac1{\sqrt{(k+1)(n-k+1)}}.
\]
中间一半的指标中，每个分母至多 \(n+1\)，而这样的项有比例为正的
\(n\) 项，所以 \(\abs{c_n}\) 有固定正下界，不趋于零。Cauchy 乘积发散。
\end{counterexample}

\section{无穷乘积的定义和 Cauchy 准则}

给定实数列 \((p_n)\)，有限部分乘积为
\[
 P_N=\prod_{n=1}^Np_n.
\]

\begin{definition}
若 \(P_N\to P\ne0\)，称无穷乘积
\(\prod_{n=1}^\infty p_n\) \term{收敛}，并以 \(P\) 为其值。
若 \(P_N\to0\)，称乘积\term{趋于零}，但按本书约定不称为收敛。
\end{definition}

非零极限的约定使有限个因子的改变只按一个非零常数缩放乘积，也使尾乘积
趋于 \(1\)。

\begin{theorem}[无穷乘积的 Cauchy 准则]
\label{thm:infinite-product-cauchy}
假设每个 \(p_n\ne0\)。无穷乘积收敛，当且仅当对任意
\(\varepsilon>0\)，存在 \(N\)，使 \(n\ge m\ge N\) 时
\[
 \abs{\prod_{k=m}^n p_k-1}<\varepsilon.
\]
\end{theorem}

\begin{proof}
若 \(P_N\to P\ne0\)，则
\[
 \prod_{k=m}^n p_k=\frac{P_n}{P_{m-1}}\to1
\]
在 \(m,n\to\infty\) 时成立。

反之，先取 \(N_0\) 使一切从 \(N_0\) 开始的有限尾乘积落在
\((1/2,3/2)\)。固定 \(m=N_0\)，可知
\(P_n/P_{N_0-1}\) 有界，因而 \((P_n)\) 最终有界。记
\(\abs{P_n}\le C\)（\(n\ge N_0\)）。给定 \(\eta>0\)，再由题设取
\(N\ge N_0\)，使 \(n>\ell\ge N\) 时
\[
 \abs{\frac{P_n}{P_\ell}-1}<\frac{\eta}{C}.
\]
于是
\[
 \abs{P_n-P_\ell}
 =\abs{P_\ell}\abs{\frac{P_n}{P_\ell}-1}<\eta.
\]
所以 \((P_n)\) 是 Cauchy 列，设其极限为 \(P\)。另一方面，对
\(n\ge m\ge N_0\) 的尾乘积有下界 \(1/2\)。先令 \(n\to\infty\)，
再固定任一 \(m\ge N_0\)，可得
\(\abs{P/P_{m-1}}\ge1/2\)，从而 \(P\ne0\)。
\end{proof}

\section{对数判别、零因子与乘积表示}

以下主要写 \(p_n=1+u_n\)。若乘积收敛，则尾因子
\[
 1+u_n=\frac{P_n}{P_{n-1}}\to1,
\]
所以 \(u_n\to0\) 是必要条件。

早期因子的作用必须与尾部判别分开。一个零因子使此后所有部分乘积都等于
零；有限个非零因子只把尾乘积乘以一个非零常数。对数判别处理的是最终
接近 \(1\) 的尾部，有限前缀的符号由其本身决定。

\begin{theorem}[正因子的对数判别]
\label{thm:infinite-product-log}
假设没有零因子，并且 \(1+u_n>0\) 最终成立。则
\[
 \prod_{n=1}^\infty(1+u_n)
\]
收敛，当且仅当从某个充分大的 \(N_0\) 起的对数级数
\[
 \sum_{n=N_0}^\infty\log(1+u_n)
\]
收敛于有限实数；乘积趋于零，当且仅当该对数级数趋于 \(-\infty\)。
\end{theorem}

\begin{proof}
取 \(N_0\) 使 \(1+u_n>0\) 对 \(n\ge N_0\) 成立。由于此前没有零因子，
有限前缀乘积 \(P_{N_0-1}\ne0\)，并且
\[
 \log\abs{P_N}=\log\abs{P_{N_0-1}}
 +\sum_{n=N_0}^N\log(1+u_n).
\]
尾部各因子为正，故 \(P_N\) 的符号最终固定。指数函数连续且把有限实数
映成正数，所以对数部分和有有限极限，当且仅当 \(P_N\) 有非零极限。
若对数部分和趋于 \(-\infty\)，则 \(\abs{P_N}\to0\)；反向也成立。
\end{proof}

\begin{corollary}
若 \(\sum\abs{u_n}<\infty\)，则
\(\prod(1+u_n)\) 在没有零因子时收敛。若 \(u_n\ge0\)，则
\[
 \prod(1+u_n)\text{ 收敛}
\quad\Longleftrightarrow\quad
\sum u_n\text{ 收敛}.
\]
\end{corollary}

\begin{proof}
充分大时 \(\abs{u_n}\le1/2\)，并有
\[
 \abs{\log(1+u_n)}\le2\abs{u_n}.
\]
故绝对收敛推出对数级数收敛。若 \(u_n\ge0\) 且趋零，则充分大时
\[
 \frac12u_n\le\log(1+u_n)\le u_n,
\]
比较判别给出等价性。
\end{proof}

条件收敛的 \(\sum u_n\) 仍可能给出收敛乘积，但需控制二次误差。
例如从 \(n=2\) 开始取 \(u_n=(-1)^n/n\)，则
\[
 \log(1+u_n)=u_n+O(u_n^2).
\]
第一项级数收敛，第二项绝对可和，所以乘积收敛到非零数。
这里使用的是：当 \(\abs u\le1/2\) 时，存在常数 \(C\) 使
\[
 \abs{\log(1+u)-u}\le C u^2.
\]
因此“\(\sum u_n\) 收敛且 \(\sum u_n^2\) 收敛”足以控制取对数产生的
累计误差；只有 \(\sum u_n\) 收敛一般还不够。

\section*{习题}
\addcontentsline{toc}{section}{习题}

\subsection*{A组　二重求和}
\begin{enumerate}[label=\textbf{\thechapter.\arabic*.}]
 \exerciseitem{double-series-counterexample-verify} 对正文反例逐项计算行和、列和、
 \(S_{M,M}\) 与 \(S_{M,M+1}\)。
 \exerciseitem{double-series-diagonal-not-enough} 构造二重数列使
 \(S_{N,N}\to0\)，但 Pringsheim 极限不存在。
 \exerciseitem{double-series-nonnegative-supremum} 从有限矩形上确界出发，
 完整证明非负二重级数两个迭代和相等。
 \exerciseitem{double-series-absolute-enumeration} 证明绝对收敛二重级数按
 \(\Npos^2\) 的任意枚举求和都得到同一值。
 \exerciseitem{double-series-product-separable} 若
 \(\sum\abs{a_m}<\infty\)、\(\sum\abs{b_n}<\infty\)，证明
 \[
  \sum_{m,n\ge1}a_mb_n
  =\left(\sum_ma_m\right)\left(\sum_nb_n\right).
 \]
\end{enumerate}

\subsection*{B组　Cauchy 乘积}
\begin{enumerate}[label=\textbf{\thechapter.\arabic*.},resume]
 \exerciseitem{cauchy-product-geometric} 直接计算两个几何级数的 Cauchy
 乘积，并核对总和公式。
 \exerciseitem{cauchy-product-partial-sum-identity} 证明
 \[
  \sum_{n=0}^Nc_n=\sum_{k=0}^Na_kB_{N-k}
 \]
 并用图示或指标区域说明缺失的是哪一块二重尾项。
 \exerciseitem{mertens-proof-tail} 把 Mertens 定理证明中的有限头部与绝对
 可和尾部估计写成完整的 \(\varepsilon\) 证明。
 \exerciseitem{cauchy-product-conditional-divergence} 对
 \(a_n=(-1)^n/\sqrt{n+1}\) 给出 \(\abs{c_n}\) 的显式统一正下界。
 \exerciseitem{cauchy-product-one-finite} 若一个级数只有有限个非零项，证明
 它与任意收敛级数的 Cauchy 乘积收敛到两和之积。
\end{enumerate}

\subsection*{C组　无穷乘积}
\begin{enumerate}[label=\textbf{\thechapter.\arabic*.},resume]
 \exerciseitem{infinite-product-telescoping} 计算
 \[
  \prod_{n=2}^\infty\left(1-\frac1{n^2}\right).
 \]
 \exerciseitem{infinite-product-zero-harmonic} 证明
 \(\prod_{n=2}^\infty(1-1/n)=0\)，并解释为何本书不称其为收敛乘积。
 \exerciseitem{infinite-product-absolute-perturbation} 证明若
 \(\sum\abs{u_n}<\infty\) 且无零因子，则改变有限多个因子只把乘积乘以
 一个非零常数。
 \exerciseitem{infinite-product-log-second-order} 若 \(u_n\to0\)、
 \(\sum u_n\) 收敛且 \(\sum u_n^2<\infty\)，证明
 \(\prod(1+u_n)\) 在无零因子时收敛。
 \exerciseitem{infinite-product-positive-divergence} 若 \(u_n\ge0\) 且
 \(\sum u_n=\infty\)，证明部分乘积趋于 \(+\infty\)。
 \optionalexerciseitem{euler-product-finite-stage}
 对有限素数集 \(P\) 和 \(s>1\)，展开有限乘积
 \[
  \prod_{p\in P}\frac1{1-p^{-s}},
 \]
 说明每个只含 \(P\) 中素因子的正整数恰出现一次。再讨论把 \(P\) 扩大时，
 这些有限乘积与 \(\sum_{n\ge1}n^{-s}\) 的关系。
\end{enumerate}
